% ============================================================================
% CML Beamer Template
% Matches the visual style used in Computational Memory Lab lectures
% (Madrid theme + navy blue + lab-branded footer).
%
% Variables to fill in (look for <<...>> markers):
%   <<TITLE>>       main talk title
%   <<SUBTITLE>>    optional subtitle (delete the \subtitle line if unused)
%   <<AUTHOR>>      presenter name(s)
%   <<DATE>>        defaults to today
%   <<SHORT_TITLE>> short version shown in footer (max ~5 words)
%
% Use \section{...} for each major section. Each section auto-updates the
% navigation dots shown in the Madrid header.
%
% For multi-paper talks, typical structure:
%   \section{Paper 1}
%   \section{Paper 2}
%   \section{Comparison}
%
% For figures, drop files into a `figures/` directory next to this .tex file
% and reference relatively: \includegraphics[width=...]{figures/myfig.pdf}
% ============================================================================

\documentclass[aspectratio=169]{beamer}

% --- Theme -------------------------------------------------------------------
\usetheme{Madrid}
\usecolortheme{default}

% CML navy blue — the shade used in the BigData_ML_Lecture reference
\definecolor{cmlblue}{RGB}{28, 53, 115}
% On Madrid, recoloring `structure` + the palette cascades cleanly to the
% title banner, the headline strip, the footer, and frame titles — without
% the title-text ending up the same color as its background. Don't override
% individual elements (title, frametitle) here unless you're willing to
% double-check that fg and bg haven't collided.
\setbeamercolor{structure}{fg=cmlblue}
\setbeamercolor{palette primary}{bg=cmlblue, fg=white}
\setbeamercolor{palette secondary}{bg=cmlblue!80, fg=white}
\setbeamercolor{palette tertiary}{bg=cmlblue!60, fg=white}

% --- Packages ---------------------------------------------------------------
\usepackage[T1]{fontenc}
\usepackage[utf8]{inputenc}
\usepackage{amsmath, amssymb, amsthm}
\usepackage{graphicx}
\usepackage{hyperref}
\hypersetup{colorlinks=true, linkcolor=cmlblue, urlcolor=cmlblue, citecolor=cmlblue!70}
\usepackage[round, authoryear]{natbib}
% Bibliography is loaded via \bibliography{} in the References frame near
% the end of the document. Natbib + bibtex works out of the box on any
% TeX Live install. Citation commands: \citet{key}, \citep{key}, \citeauthor{key}.

% --- Notes configuration ----------------------------------------------------
% Two compile modes:
%   default: talk-facing (notes hidden)
%   uncomment the show-only-notes line + recompile for a notes-only handout
\setbeameroption{hide notes}
% \setbeameroption{show only notes}

% --- Metadata ---------------------------------------------------------------
\title[<<SHORT_TITLE>>]{<<TITLE>>}
\subtitle{<<SUBTITLE>>}
\author[<<AUTHOR>>]{<<AUTHOR>>}
\institute[Penn CML]{Computational Memory Lab \\ University of Pennsylvania}
\date{<<DATE>>}

% ============================================================================
\begin{document}

% --- Title slide -----------------------------------------------------------
\begin{frame}[plain]
  \titlepage
\end{frame}

% --- Outline ---------------------------------------------------------------
\begin{frame}{Outline}
  \tableofcontents
\end{frame}

% ============================================================================
% Sections begin here. Replace the templates below with actual content.
% ============================================================================

\section{Motivation}

\begin{frame}{Why this matters}
  \begin{itemize}
    \item Key point 1 framed for the audience
    \item Key point 2 — why the existing answer is insufficient
    \item What this talk adds
  \end{itemize}
  \note{Set the stakes. Keep to 45 seconds. The audience should
        know by the end of this slide what question the talk answers.}
\end{frame}

\section{Background}

\begin{frame}{Prior work}
  \begin{itemize}
    \item Existing model / result that everyone knows
    \item Specific gap or puzzle we're addressing
  \end{itemize}
  \note{}
\end{frame}

\section{Methods}

\begin{frame}{Experimental design}
  % Figure placeholder — replace with actual figure from bundle
  % \includegraphics[width=0.9\textwidth]{figures/design.pdf}
  \begin{itemize}
    \item Participants: N = ?
    \item Task: ?
    \item Key analyses: ?
  \end{itemize}
  \note{}
\end{frame}

\section{Results}

\begin{frame}{Headline finding}
  \begin{columns}
    \column{0.6\textwidth}
      % \includegraphics[width=\textwidth]{figures/fig_main.pdf}
    \column{0.4\textwidth}
      \begin{itemize}
        \item What the figure shows
        \item Why it's surprising / confirmatory
      \end{itemize}
  \end{columns}
  \note{The punchline slide. Don't talk to the bullets — talk to the
        figure. The bullets are just for the audience's eyes.}
\end{frame}

\section{Discussion}

\begin{frame}{Implications}
  \begin{itemize}
    \item What this changes about our understanding
    \item Open question for future work
  \end{itemize}
  \note{}
\end{frame}

% --- Acknowledgments -------------------------------------------------------
\begin{frame}{Acknowledgments}
  \begin{itemize}
    \item Collaborators, funding sources
  \end{itemize}
\end{frame}

% --- References -----------------------------------------------------------
\begin{frame}[allowframebreaks]{References}
  \bibliographystyle{apalike}
  \bibliography{references}
\end{frame}

\end{document}
